\section{Introduction}
Energetic molecular crystals such as CL-20 and HMX remain a difficult regime for atomistic simulation because compression, intermolecular rearrangement, bond breaking, and charge redistribution unfold on comparable femto--picosecond time scales. In that setting, a potential must remain reliable not only near equilibrium, but also for distorted configurations in which polarization and changing molecular environments can alter both relative energetics and subsequent trajectory evolution \cite{valeeswaran2023charge,menikoff2016shock,bolton2008molecular}.

Recent machine-learned interatomic potentials have substantially improved the accuracy and scale of atomistic simulation, particularly through equivariant architectures that encode geometric symmetry directly \cite{batatia2022mace,batzner20223,musaelian2023learning,liao2023equiformer,batatia2023foundation}. At the same time, recent reviews and transferability studies have made clear that accuracy, transfer, and execution cost remain tightly coupled design constraints rather than independent objectives \cite{musil2021physics,unke2021machine,goedecker2023perspective,jacobs2025practical,huang2025cross}.

For reactive molecular crystals, a central ambiguity therefore remains. When a finite-cutoff model fails under compression or along nonequilibrium trajectories, the missing ingredient may be a richer short-range representation, but it may equally be an explicit treatment of environment-dependent electrostatic response. That distinction is especially important for energetic materials, where polarization and charge transfer have long been implicated in sensitivity, decomposition, and shock response \cite{valeeswaran2023charge,van2001reaxff}.

Charge-aware machine-learning potentials offer a natural route toward that missing physics. Charge-equilibration ideas originate in the QEq formalism of Rapp\'e and Goddard \cite{rappe1991charge} and now appear in several atomistic-learning frameworks that couple learned local features to environment-dependent charges or electrostatic observables \cite{gastegger2017machine,unke2019physnet,ko2021fourth,fuchs2025celli}. Contemporary materials studies have also begun to test such ideas directly under non-equilibrium conditions. The question is therefore no longer whether long-range response can be embedded in an MLIP, but how that response changes structural stability once trajectories leave the near-equilibrium manifold \cite{jung2026charge}.

What remains unresolved is whether explicit charge self-consistency improves the observables that matter most for energetic-material modelling. A better electrostatic description may sharpen static energy ranking without stabilizing the nonequilibrium quantities that govern bond survival, fragment integrity, and trajectory degradation. The relevant question is therefore not simply whether charge-aware models lower a benchmark error, but whether they alter the failure structure of trajectories that are already leaving the near-equilibrium manifold.

We address that question with CSC-MACE, a charge-self-consistent extension of MACE in which equivariant atomic features parameterize a differentiable charge-equilibration contribution to energies, forces, and stresses within one periodic total-energy model. The model is evaluated across a layered benchmark ladder comprising a public reference system (OpenREACT), two public surrogate systems (Transition1x and SPICE), and a final CL-20/HMX benchmark. This structure separates generic public-data behavior from the chemically decisive target and allows two linked questions to be tested directly: whether explicit self-consistent electrostatic response improves static fitting on a real energetic-material benchmark, and whether any such gain propagates into a reproducible nonequilibrium hierarchy under shock-motivated rollout protocols.

On the final CL-20/HMX benchmark, charge-aware models improve the static energy hierarchy, and the self-consistent branch remains best on energy error under four held-out views: the original contiguous slice, the cleaner material-balanced split-v2 evaluation, the stronger order-recovered block-holdout split-v3 evaluation, and the stronger recovered temperature-gap split-v4 evaluation.

The dynamical evidence points elsewhere. Across stress protocols, frame-indexed survival analyses, and train-matched fragile-frame comparisons, degradation remains dominated primarily by starting-frame choice rather than by model family, with frame-driven variation exceeding family-driven variation by about $8\times$ for the first integrity-crossing step and about $184\times$ for the first drift-threshold step in the focused final-target comparison. A denser failure-window analysis nevertheless identifies a narrower positive signal: explicit self-consistency does not create a broad onset advantage, but in a narrow fragile window around frames $64$--$72$ it improves endpoint structural survival, with the clearest local gain at frame $72$.

The paper therefore establishes a more precise result: explicit electrostatic response is physically active and beneficial for static energetic-material fitting, while its transfer to nonequilibrium stability remains strongly bounded under the present rollout protocols.
